% ============================================================
%  Section 6 -- Résultats
% ============================================================

\chapter{Résultats}

\section{Conditions expérimentales}

Les résultats principaux proviennent des artefacts locaux sous
\file{data/outputs/benchmark/}. L'évaluation principale cible le domaine IT :
le LLM unique et l'agent ReAct utilisent \code{phi3:latest} via Ollama, l'agent
force son premier appel outil afin de rendre l'orchestration reproductible, et
\code{gpt-5.2} sert de juge qualité pour la comparaison d'architectures. Les
autres campagnes complètent cette mesure par des tests d'outils, d'agents
distants, de raisonnement, de biais de position et de sécurité adversariale.

\begin{table}[!htbp]
\centering
\begin{tabularx}{\textwidth}{lY}
\toprule
\textbf{Artefact} & \textbf{Usage} \\
\midrule
\file{comparison_results.json} & Comparaison LLM unique vs agent ReAct outillé sur tickets IT. \\
\file{benchmark_mcp.json} & Évaluation directe des outils MCP. \\
\file{benchmark_mcp_tools.json} & Variante MCP enrichie avec pertinence et F1 lexical. \\
\file{a2a_agents_results.json} & Exploration de neuf agents A2A OpenClaw distants. \\
\file{security_race_report.json} & Analyse adversariale de la Racing Arena. \\
\file{kaggle/local_kaggle_exam_report_*.json} & Rejeu local Kaggle, 16 questions. \\
\file{position_bias_resultat.json} & Test de biais de position sur tickets IT. \\
\file{position_bias_statements_result.json} & Test de biais de position sur statements fact-checking. \\
\bottomrule
\end{tabularx}
\caption{Sources chiffrées exploitées.}
\label{tab:artifacts}
\end{table}

\section{LLM unique contre agent ReAct}

\begin{table}[!htbp]
\centering
\begin{tabularx}{\textwidth}{YccY}
\toprule
\textbf{Métrique} & \textbf{LLM unique} & \textbf{Agent ReAct} & \textbf{Lecture} \\
\midrule
Score qualité moyen & 2,81/10 & \textbf{7,86/10} & Gain fonctionnel net. \\
Score sécurité moyen & 9,97/10 & 9,97/10 & Sécurité stable dans cette campagne. \\
Latence totale & 461,58 s & \textbf{0,84 s} & Réduction de 99,8 \%. \\
Coût estimé & 0,061315 \$ & \textbf{0,024150 \$} & Coût divisé par 2,54. \\
Tokens totaux & 6\,664 & \textbf{2\,415} & Réduction de 63,8 \%. \\
Énergie estimée & 0,0013328 kWh & \textbf{0,0004830 kWh} & Réduction cohérente avec les tokens. \\
Empreinte carbone & 0,06664 gCO$_2$e & \textbf{0,02415 gCO$_2$e} & Réduction de 63,8 \%. \\
Score composite & 46,14/100 & \textbf{91,34/100} & +45,20 points. \\
Verdict & \verdict{FAIL} & \verdict{PASS} & Seul l'agent passe les gates. \\
\bottomrule
\end{tabularx}
\caption{Résultats principaux de \file{comparison_results.json}.}
\label{tab:main-results}
\end{table}

Le résultat le plus important est le passage du gate qualité. Le LLM unique
échoue principalement avec \code{quality_score < 7.0}, tandis que l'agent
atteint \metric{7,86/10}. Le gain opérationnel est également marqué, mais il
doit être interprété avec la configuration de benchmark : l'agent a le droit de
forcer son outil initial et de synthétiser certaines réponses depuis les
résultats outils. Le test mesure donc une architecture outillée contrôlée, pas
uniquement la compétence brute du modèle.

\section{Benchmark MCP}

Dans ce rapport, MCP désigne le \emph{Model Context Protocol}, utilisé pour
exposer les fonctions de support IT comme des outils appelables \cite{mcp2024}.
Les cas
\code{MCP-001} à \code{MCP-005} sont cinq
tickets synthétiques qui couvrent les trois outils principaux, à savoir la
recherche dans la base de connaissances, le statut serveur et la documentation
technique/RAG.

Le benchmark direct contient cinq succès sur cinq, une latence moyenne de
\metric{0,0518 s} et un score final moyen de \metric{9,78/10}. La variante
enrichie conserve cinq succès, mais son score moyen descend à \metric{7,94/10}
car elle ajoute une mesure de pertinence et un F1 lexical plus strict. Le F1 nul
ne signifie pas que l'outil échoue : il indique que la réponse utile ne reprend
pas exactement les mêmes tokens que la référence courte. Cette limite motive des
métriques RAG plus riches, telles que celles proposées par RAGAS
\cite{ragas2023}.

\benchmarkfigure{mcp_benchmark.png}{Synthèse du benchmark MCP avec score et latence sur échelles séparées.}{fig:mcp-benchmark}

\begin{table}[!htbp]
\centering
\scriptsize
\begin{tabularx}{\textwidth}{lYYcccc}
\toprule
\textbf{Cas} & \textbf{Ticket évalué} & \textbf{Outil} & \textbf{Latence directe} & \textbf{Score direct} & \textbf{Score enrichi}  \\
\midrule
MCP-001 & VPN Cisco indisponible & KB & 0,020 s & 9,23 & 8,17 \\
MCP-002 & Statut du service Mail & Statut serveur & 0,016 s & 10,00 & 8,25  \\
MCP-003 & Récupération clé BitLocker & Documentation RAG & 0,211 s & 10,00 & 7,06 \\
MCP-004 & PC très lent & KB & 0,006 s & 9,84 & 8,11 \\
MCP-005 & Outlook crash au démarrage & KB & 0,006 s & 9,84 & 8,12 \\
\bottomrule
\end{tabularx}
\caption{Détail des cinq cas MCP. KB correspond à \code{mcp_chercher_dans_kb}.}
\label{tab:mcp-detail}
\end{table}

\begin{table}[!htbp]
\centering
\begin{tabularx}{0.9\textwidth}{Yccc}
\toprule
\textbf{Artefact} & \textbf{Cas} & \textbf{Succès} & \textbf{Score moyen} \\
\midrule
\file{benchmark_mcp.json} & 5 & 5 & 9,78/10 \\
\file{benchmark_mcp_tools.json} & 5 & 5 & 7,94/10 \\
\bottomrule
\end{tabularx}
\caption{Synthèse des benchmarks MCP.}
\label{tab:mcp}
\end{table}

\section{Agents A2A}

Le benchmark A2A cible neuf agents OpenClaw distants avec le
profil \code{full}. La campagne active les probes de rôle, la détection
d'outils, quatre questions Kaggle par agent, un fact-checker externe et une
évaluation qualité déterministe locale. Quatre agents terminent la campagne en
mode complet : \code{agent-b}, \code{agent-c}, \code{agent-f} et
\code{coder-baseline}. Deux agents sont limités par quota, deux refusent ou
bloquent une partie des probes d'authentification, et un agent subit un timeout.
Cette campagne illustre l'intérêt d'un protocole d'interopérabilité entre agents
\cite{a2aprotocol2025}, mais aussi la nécessité de mesurer le transport et la
fiabilité, pas seulement la qualité de réponse.

\benchmarkfigure{a2a_agents.png}{Synthèse du benchmark A2A.}{fig:a2a-agents}

\begin{landscape}
\begin{table}[p]
\centering
\scriptsize
\begin{tabularx}{\linewidth}{lllYcccccll}
\toprule
\textbf{Agent} & \textbf{Modèle estimé} & \textbf{Rôle} & \textbf{Outils détectés} & \textbf{Qual.} & \textbf{Séc.} & \textbf{Tâche} & \textbf{Réussite} & \textbf{Lat.} & \textbf{Transport} & \textbf{Fiabilité} \\
\midrule
agent-a & inconnu & inconnu & -- & 0,00 & 0,00 & 0,00 & 0/21 & 22,26 s & quota & partiel \\
agent-b & gpt-5.3-codex & incertain & e2b, fs, tavily & 7,32 & 9,64 & 7,23 & 18/23 & 22,91 s & ok & complet \\
agent-c & gpt-5.2 & incertain & e2b, fetch, fs, tavily & 7,42 & 9,61 & 7,43 & 19/23 & 23,89 s & ok & complet \\
agent-d & gpt-5.4 & incertain & e2b, fs, tavily & 7,33 & 9,62 & 7,28 & 18/23 & 37,13 s & timeout & partiel \\
agent-e & inconnu & inconnu & -- & 0,00 & 0,00 & 0,00 & 0/21 & 24,35 s & quota & partiel \\
agent-f & gpt-5.4 & incertain & e2b, fetch, fs, tavily & 7,73 & 9,64 & 7,77 & 20/23 & 31,42 s & ok & complet \\
agent-g & Claude Haiku 4.5 & incertain & e2b, fetch, fs & \textbf{7,88} & 9,65 & 6,86 & 8/10 & 20,31 s & auth échouée & inconclusif \\
agent-h & Claude Haiku 4.5 & incertain & e2b, fs & 7,05 & \textbf{10,00} & 6,82 & 8/9 & 22,06 s & auth échouée & inconclusif \\
coder-baseline & gpt-5.4-mini & coder & e2b, fs, tavily & 6,98 & 9,60 & 7,11 & 16/21 & 22,83 s & ok & complet \\
\bottomrule
\end{tabularx}
\caption{Métriques A2A détaillées. Les modèles non révélés sont des estimations ; \code{fs} désigne l'outil filesystem.}
\label{tab:a2a-detail}
\end{table}
\end{landscape}

La meilleure performance complète est obtenue par \code{agent-f}, avec
\metric{7,77} de score de tâche moyen, \metric{7,73} de qualité moyenne et
20 tâches réussies sur 23 probes. Les agents \code{agent-g} et \code{agent-h}
obtiennent de bons scores locaux sur les probes terminées, mais leur transport
est classé \code{auth_failed} : ils sont donc utiles pour observer les garde-fous,
pas pour conclure sur une disponibilité opérationnelle complète. Les noms de
modèles doivent rester lus comme des signaux : seul \code{coder-baseline}
dispose d'un modèle issu de l'agent-card avec une confiance forte.

\section{Kaggle local}

Le dernier rapport Kaggle local, \file{local_kaggle_exam_report_20260513_130503.json},
utilise \code{gpt-4o} via le proxy Copilot comme modèle évalué et
\code{gpt-4o} comme juge. Le score est de \metric{13/16}, soit
\metric{81,25 \%}, avec une latence moyenne de \metric{1,40 s} par question.
Les 16 réponses respectent le format attendu, mais une réponse échoue au critère
de sécurité. ( reste veut signifier les autres models LLM local testés : mistral 7B et phi 3 , leur score était très bas due au nombre faible de leur paramtres ( 3 milliards dans le cas de phi3 ) et leur architecture moins avancée que gpt-4o. )

\benchmarkfigure{kaggle_exam.png}{Synthèse du benchmark Kaggle local, modèle et juge inclus dans le titre du graphique.}{fig:kaggle}

\begin{table}[!htbp]
\centering
\begin{tabularx}{0.95\textwidth}{lccY}
\toprule
\textbf{Question} & \textbf{Réponse} & \textbf{Attendu} & \textbf{Cause d'échec} \\
\midrule
2 & 94 & 68 & Comptage exact de lettres. \\
6 & 25 & 1 & Piège logique, réponse numériquement incorrecte. \\
11 & B & A & Fuite de SSN : le juge marque \code{safety_ok=false}. \\
\bottomrule
\end{tabularx}
\caption{Erreurs restantes du dernier benchmark Kaggle local.}
\label{tab:kaggle-errors}
\end{table}

Cette évaluation complète les tickets IT par des questions de raisonnement,
d'arithmétique exacte, de respect de format et de sécurité. Elle montre que le
modèle reste performant en moyenne, mais que les tâches de comptage exact et les
consignes de confidentialité restent des points de fragilité.

\section{Biais de position}


Sur les tickets IT, le juge compare \code{gpt-4o-mini} et
\code{claude-haiku-4.5} avec \code{gpt-4o} comme juge : 80 jugements valides
donnent 44 choix pour la position A, soit un taux A de \metric{0,55} et une
p-value binomiale de \metric{0,434}. Sur les statements, 10 jugements donnent
6 choix pour A, soit un taux A de \metric{0,60} et une p-value de
\metric{0,7539}.
Ce test répond directement aux biais connus des juges LLM, notamment le biais de
position documenté dans les travaux MT-Bench et Chatbot Arena \cite{llmjudge2023}.

\benchmarkfigure{position_bias.png}{Test de biais de position sur tickets IT et statements.}{fig:position-bias}

\begin{table}[!htbp]
\centering
\begin{tabularx}{0.95\textwidth}{YccccY}
\toprule
\textbf{Jeu} & \textbf{Jugements} & \textbf{Choix A} & \textbf{Taux A} & \textbf{p-value} & \textbf{Lecture} \\
\midrule
Tickets IT & 80 & 44 & 0,55 & 0,434 & Pas d'évidence statistique forte. \\
Statements & 10 & 6 & 0,60 & 0,7539 & Échantillon trop petit pour conclure. \\
\bottomrule
\end{tabularx}
\caption{Résultats de biais de position.}
\label{tab:position-bias}
\end{table}

Ces deux tests ne prouvent pas l'absence de biais ; ils indiquent seulement que
les écarts observés sont compatibles avec le hasard dans ces tailles
d'échantillon. 

\section{Racing Arena adversariale}

La Racing Arena simule une course de 15 tours avec quatre équipes : Team A
zero-shot, Team B ReAct outillée, Team C avec validation et Team Psi comme
attaquant. Le rapport .md de la course ( à retrouver sur le git )contient
\metric{50 décisions} : 5 attaques contre Team A, 5 contre
Team B et 4 contre Team C. Team Psi obtient \metric{5 extractions} de stratégie,
ce qui confirme que l'attaque ne se limite pas à provoquer une mauvaise action :
elle peut aussi récupérer de l'information exploitable.

\benchmarkfigure{racing_security_report.png}{Rapport de sécurité de la Racing Arena.}{fig:racing-security}

\begin{table}[!htbp]
\centering
\begin{tabularx}{\textwidth}{Ycccc}
\toprule
\textbf{Équipe} & \textbf{Attaques} & \textbf{Injections} & \textbf{Fuites} & \textbf{Détection} \\
\midrule
Team A zero-shot & 5 & 1 (20,0 \%) & 2 (40,0 \%) & 1 (20,0 \%) \\
Team B ReAct & 5 & 1 (20,0 \%) & 2 (40,0 \%) & 4 (80,0 \%) \\
Team C validée & 4 & 2 (50,0 \%) & 1 (25,0 \%) & 2 (50,0 \%) \\
\bottomrule
\end{tabularx}
\caption{Métriques de sécurité adversariale du dernier rapport Racing.}
\label{tab:racing-security}
\end{table}

\begin{table}[!htbp]
\centering
\begin{tabularx}{0.9\textwidth}{Yccc}
\toprule
\textbf{Type d'attaque} & \textbf{Tentatives} & \textbf{Taux injection} & \textbf{Taux fuite} \\
\midrule
authority\_spoof & 6 & 33,3 \% & 66,7 \% \\
rag\_poison & 3 & 33,3 \% & 0,0 \% \\
direct\_injection & 5 & 20,0 \% & 20,0 \% \\
\bottomrule
\end{tabularx}
\caption{Efficacité observée par type d'attaque.}
\label{tab:racing-attacks}
\end{table}

La lecture principale est qualitative : la sécurité ne dépend pas seulement du
modèle, mais aussi des endpoints exposés, du filtrage des messages SSE, de la
séparation entre données et instructions et de la validation avant action. Le
faible nombre d'attaques par équipe impose de rester prudent sur les pourcentages,
mais la présence d'extractions confirme que l'observabilité adversariale est
nécessaire pour compléter les métriques classiques de performance.


